% Iteration 8 change manifest as two structured cards: chg-1 hardens the
% publish-state shell guard, chg-2 makes execution-risk warnings salient on the
% next model turn and adds contract-aware heuristics. Both ship together at the
% iteration-7 boundary as a deliberate keep-and-improve on the architecture
% installed by iterations 5 and 6.
\begin{figure}[t]
    \centering
    \scriptsize
    \begin{tcbraster}[raster columns=2,
                      raster equal height,
                      raster column skip=4pt,
                      raster row skip=0pt,
                      raster left skip=0pt,
                      raster right skip=0pt,
                      boxrule=0.6pt,
                      left=4pt,right=4pt,top=3pt,bottom=3pt,
                      coltitle=white]
        \begin{tcolorbox}[colback=aheSky!18,colframe=aheNavy,colbacktitle=aheNavy,
                          title={\scriptsize\textbf{chg-1, iteration 8, commit \texttt{ca35f53}, level: tool implementation}}]
            \textbf{\textcolor{aheNavy}{Files}} \\
            \rule{\linewidth}{0.3pt}\\[1pt]
            $\bullet$ \texttt{tools/shell\_tools/run\_shell\_command.py} \\[3pt]
            \rule{\linewidth}{0.3pt}\\
            \textbf{\textcolor{aheNavy}{What changed}} \\
            \rule{\linewidth}{0.3pt}\\[1pt]
            Upgraded two soft reasons to hard blocks. Deletion of any non-\texttt{/tmp} protected output is now a hard block. Reset of any non-\texttt{/tmp} protected root to an empty state is also a hard block. The \texttt{ALLOW\_POST\_SUCCESS\_RESET} token still exists for other classes of post-success interlock but can no longer wipe verified live deliverables or empty live roots. \\[3pt]
            \rule{\linewidth}{0.3pt}\\
            \textbf{\textcolor{aheNavy}{Failure pattern fixed}} \\
            \rule{\linewidth}{0.3pt}\\[1pt]
            Agent attached the override token to delete \texttt{/git/www/hello.html} and reset \texttt{/git/server/refs/heads/master} after a successful deployment check, ``returning to a clean repo''; verifier then 404s. \\[3pt]
            \rule{\linewidth}{0.3pt}\\
            \textbf{\textcolor{aheNavy}{Predicted fixes}} \\
            \rule{\linewidth}{0.3pt}\\[1pt]
            2 tasks. Examples: \texttt{configure-git-webserver}, \texttt{git-multibranch}.
        \end{tcolorbox}
        \begin{tcolorbox}[colback=aheTeal!12,colframe=aheBlue,colbacktitle=aheBlue,
                          title={\scriptsize\textbf{chg-2, iteration 8, commit \texttt{a4a4a29}, level: middleware}}]
            \textbf{\textcolor{aheBlue}{Files}} \\
            \rule{\linewidth}{0.3pt}\\[1pt]
            $\bullet$ \texttt{workspace/middleware/execution\_risk\_hints.py} \\[3pt]
            \rule{\linewidth}{0.3pt}\\
            \textbf{\textcolor{aheBlue}{What changed}} \\
            \rule{\linewidth}{0.3pt}\\[1pt]
            Added a \texttt{BeforeModelHook} that promotes any execution-risk note emitted on the previous step into a FRAMEWORK reminder visible at the top of the next model turn, so warnings enter the reasoning context rather than trail after the tool output. Added one-time per-task contract inference for clean-layout or single-file delivery contracts and official-wrapper or named-revision contracts. Added two new after-tool heuristics: a warning when the agent compiles or builds inside a clean-layout live tree, and a warning when the contract names an official wrapper but the command uses a raw \texttt{SentenceTransformer} or \texttt{AutoModel} style API instead. \\[3pt]
            \rule{\linewidth}{0.3pt}\\
            \textbf{\textcolor{aheBlue}{Failure pattern fixed}} \\
            \rule{\linewidth}{0.3pt}\\[1pt]
            Iteration-6 middleware emitted the right warnings but only into tool output; the agent often made the publish/stop decision on the next model turn and ignored them. Salience promotion plus contract-aware heuristics close this gap. \\[3pt]
            \rule{\linewidth}{0.3pt}\\
            \textbf{\textcolor{aheBlue}{Predicted fixes}} \\
            \rule{\linewidth}{0.3pt}\\[1pt]
            4 tasks. Examples: \texttt{polyglot-c-py}, \texttt{polyglot-rust-c}, \texttt{mteb-retrieve}, \texttt{pytorch-model-recovery}.
        \end{tcolorbox}
    \end{tcbraster}
    \caption{Two change-manifest entries written together at the iteration-7 boundary and shipped as the iteration-8 harness. \texttt{chg-1} hardens the existing publish-state shell guard so that the override token can no longer wipe verified live deliverables. \texttt{chg-2} makes execution-risk warnings impossible to overlook at the next model turn and adds two contract-aware heuristics. Both are deliberately scoped: \texttt{chg-1} prevents the destructive command itself, \texttt{chg-2} fixes the salience gap of the iteration-6 middleware.}
    \label{fig:manifest-hardblock}
\end{figure}
